
\medskip

L'image ci-dessous représente la position obtenue au déclenchement du bloc départ
d'un programme de jeu.

\medskip

\parbox{0.55\linewidth}{\psset{unit=0.5cm}
\begin{pspicture*}(-6,-4)(6,4)
\multido{\n=-6+1}{13}
{
\multido{\na=-4+1}{9}{\psdots[dotscale=0.5](\n,\na)}
}
\psaxes[Dx=10,Dy=10]{->}(0,0)(-6,-4)(6,4)
\uput[dl](0,0){O}
\psdots[dotscale=2](4,3)
\rput(-3,-2){\includegraphics[width=4mm]{chat}}
\end{pspicture*}}
\hfill
\parbox{0.43\linewidth}{L'arrière-plan est constitué de points
espacés de 40 unités.

Dans cette position, le chat a pour
coordonnées $(- 120~;~-80)$.

\textbf{Le but du jeu est de positionner le
chat sur la balle.}}

\medskip

\begin{enumerate}
\item Quelles sont les coordonnées du centre de la balle représentée dans cette position?
\item Dans cette question, le chat est dans la position obtenue au déclenchement du bloc
départ.

Voici le script du lutin \og chat \fg{} qui se déplace.

\parbox{0.43\linewidth}{\textbf{a.} Expliquez pourquoi le chat ne revient pas à sa position
de départ si le joueur appuie sur la touche $\to$ puis sur la touche $\gets$.

\textbf{b.} Le joueur appuie sur la succession de touches
suivante : $\to$ $\to$  $\uparrow$ $\gets$ $\downarrow$.

Quelles sont les coordonnées $x$ et $y$ du chat après ce déplacement ?

\textbf{c.} Parmi les propositions de succession de touches ci-dessous, laquelle permet au
chat d'atteindre la balle ?} \hfill 
\parbox{0.51\linewidth}{\scriptsize \begin{scratch}
\blockinit{Quand \greenflag est cliqué}
\blockmove{Départ}
\end{scratch}

\begin{scratch}
\blockinit{Quand flèche gauche est cliqué}
\blockvariable{ajouter \ovalnum{- 40} à \selectmenu{x}}
\end{scratch}
\begin{scratch}
\blockinit{Quand flèche droite  est cliqué}
\blockvariable{ajouter \ovalnum{80} à \selectmenu{x}}
\end{scratch}

\medskip

\begin{scratch}
\blockinit{Quand flèche haut  est cliqué}
\blockvariable{ajouter \ovalnum{80} à \selectmenu{y}}
\end{scratch}
\begin{scratch}
\blockinit{Quand flèche bas  est cliqué}
\blockvariable{ajouter \ovalnum{$- 40$} à \selectmenu{y}}
\end{scratch}

\medskip

\begin{scratch}
\blockinit{Quand \selectmenu{n'importe quoi}  est cliqué}
{
\blockifelse{si \selectmenu{Balle} touchée alors} 
{
\blocklook{dire \txtbox{Je t'ai attrapé} pendant \ovalnum{2} secondes}
\blockmove{Départ}}
{}
}
\end{scratch}
}

\medskip

\begin{tabularx}{\linewidth}{|*{3}{>{\centering \arraybackslash}X|}}\hline
Déplacement 1 &Déplacement 2 &Déplacement 3\\ \hline
$\to \to \to \to \to \to \to \uparrow\uparrow\uparrow\uparrow\uparrow$&$\to\to\to \uparrow\uparrow\uparrow \to \downarrow \gets$&$\uparrow \to \uparrow \to \uparrow \to \to \downarrow \downarrow$\\ \hline
\end{tabularx}

\item Que se passe-t-il quand le chat atteint la balle ?
\end{enumerate}


\vspace{0,5cm}

